\documentclass[11pt,a4paper]{article}

% --- кодировка и язык ---
\usepackage{fontspec}
\usepackage{polyglossia}
\setdefaultlanguage{russian}
\setotherlanguage{english}

\setmainfont{DejaVu Serif}
\setsansfont{DejaVu Sans}
\setmonofont{DejaVu Sans Mono}

% --- разметка и типографика ---
\usepackage[margin=2.5cm]{geometry}
\usepackage{setspace}
\usepackage{csquotes}
\usepackage{enumitem}
\usepackage{xcolor}
\usepackage{booktabs}
\usepackage{graphicx}
\usepackage{amsmath}
\usepackage{hyperref}
\hypersetup{
  colorlinks=true,
  linkcolor=black,
  citecolor=blue!50!black,
  urlcolor=blue!60!black
}

% --- библиография ---
\usepackage[
  backend=biber,
  style=numeric,
  sorting=nyt,
  maxbibnames=99
]{biblatex}
\addbibresource{refs.bib}

% --- заглушки TODO ---
\newcommand{\todo}[1]{\textcolor{red!70!black}{\textbf{[TODO:} #1\textbf{]}}}

\title{%
  Разработка LLM-агентов для активного исследования\\
  частично известного северного мира с самообучением через метаагента%
}
\author{%
  Скороходов~М., Латалин~В., Молошников~И.А., Сбоев А.Г.%
  \and
  \todo{соавторы}%
}
\date{\todo{дата}}

\begin{document}
\maketitle

\begin{abstract}
В работе изучается поведение LLM-агентов в задачах активного исследования частично известного мира с текстовой (\emph{oracle}) перцепцией.
Среда моделирует арктическую (северную) станцию: динамическую проходимость
(снег, заносы, лёд), ограниченную видимость и ресурсные ограничения по заряду.
Над роботом-агентом (ReAct) надстраивается метаагент, совершенствующий промпты, скиллы и планировщик робота. В работе измеряется верхняя граница планирования при идеальном восприятии и
сравнивается LLM с детерминированными базовыми моделями по SR, SPL, покрытию и
надёжности \(\mathrm{pass}^k\).
Перенос на реального робота и физически точная модель Арктики не заявляются.
\end{abstract}

% \tableofcontents
% \newpage

% =====================================================================
\section{Введение}
\label{sec:intro}

Пространственный интеллект и активное исследование частично наблюдаемой среды~
--- один из ключевых нерешённых вопросов для foundation-моделей.
Работа Theory of Space~\cite{theoryofspace2026} показывает, что современные
LLM/VLM плохо ведут себя как активные исследователи: наблюдается
\emph{Active-Passive Gap}, неэффективное блуждание по сравнению с rule-based
программами и \emph{Belief Inertia} (модель не обновляет устаревшие представления
при противоречащих данных).

Северная/арктическая станция~--- практически значимый и малоизученный домен
для автономных роботов: динамически меняющаяся проходимость, ограниченная
видимость, ускоренный разряд аккумулятора на холоде.
Полевые исследования подтверждают сложность домена~\cite{baril2022snow}.
Готовые симуляторы слабо покрывают эту семантику, что оправдывает
создание собственной среды.

Самосовершенствующиеся агенты (self-improving / meta-agents), пишущие себе
инструменты и оптимизирующие промпты~\cite{wang2023voyager,shinn2023reflexion},
активно развиваются, но почти не изучены в контексте изменяющегося физического
мира; при этом у них есть задокументированные риски reward hacking и
misevolution~\cite{selfimprovingrisks2025,misevolution2025,rewardhacking2025}.

\paragraph{Научная рамка.}
Проект позиционируется \textbf{не} как «правдоподобный арктический симулятор
с самоулучшающимся роботом», а как лёгкий ранний прототип для изучения
LLM-агентов с явными ограничениями и абляциями.
Ключевые оговорки:
\begin{itemize}[leftmargin=*]
  \item мы измеряем \textbf{верхнюю границу} планирования при идеальном восприятии
        (oracle perception), а не способность к восприятию/навигации в реальном мире;
  \item «северная специфика» задаёт семантику и динамику задач, а не физически
        точную модель Арктики;
  \item претензии ограничены симуляцией; sim-to-real не заявляется.
\end{itemize}

\paragraph{Вклад.}
\begin{enumerate}[leftmargin=*]
  \item Северный прототип активного исследования частично известного мира
        с динамикой среды и текстовыми описаниями объектов вместо SLAM.
  \item Абляционный режим со скрытыми координатами: агент строит и ревизует
        когнитивную карту только из текстовых наблюдений
        (Construct/Revise~\cite{theoryofspace2026}).
  \item Самообучающийся метаагент с гейтом верификации: правки промптов,
        скиллов и планировщика принимаются только после регрессии на held-out.
  \item Систематическое сравнение LLM (облачные vs локальные \(\le 30\mathrm{B}\))
        с детерминированными baseline'ами по SR, SPL, покрытию и \(\mathrm{pass}^k\).
\end{enumerate}

% =====================================================================
\section{Обзор связанных работ}
\label{sec:related}

\paragraph{Агентные каркасы.}
ReAct~\cite{yao2022react} задаёт петлю «рассуждение~--- действие».
Reflexion~\cite{shinn2023reflexion} и Self-Refine~\cite{madaan2023selfrefine}
добавляют вербальную обратную связь.
Voyager~\cite{wang2023voyager} демонстрирует библиотеку переиспользуемых скиллов;
CodeAct~\cite{wang2024codeact} и ADAS~\cite{hu2024adas}~--- генерацию инструментов
и автодизайн агентов.
Promptbreeder~\cite{fernando2023promptbreeder} и DSPy~\cite{khattab2023dspy}
оптимизируют промпты.

\paragraph{Воплощённые среды и grounding.}
ALFWorld~\cite{shridhar2021alfworld} совмещает текстовую и embodied-среду;
SayCan~\cite{ahn2022saycan} grounding'ует язык в аффордансы робота;
Generative Agents~\cite{park2023generativeagents} исследуют симулякры поведения.
Ближайший прецедент «зачем свой мир»~--- ALFWorld, но у нас другая динамика
среды (северная семантика) и oracle-перцепция как сознательное упрощение.

\paragraph{Пространственная вера.}
Theory of Space~\cite{theoryofspace2026} вводит cognitive-map probing,
Active-Passive Gap и Belief Inertia~--- центральные понятия для нашей абляции
со скрытыми координатами.

\paragraph{Оценка и воспроизводимость.}
SPL~\cite{anderson2018spl}~--- стандарт для embodied-навигации;
\(\mathrm{pass}^k\)~\cite{taubench2024} отражает надёжность, а не только среднее.
Масштаб сцен обсуждается в ProcTHOR~\cite{deitke2022procthor} и
EmbodiedBench~\cite{embodiedbench2025}.
Воспроизводимость~--- по Henderson et~al.~\cite{henderson2018deeprl}
и чек-листу Pineau~\cite{pineau2018reproducibility}.
Сравнение малых и frontier-моделей опирается на AgentBench~\cite{liu2023agentbench}
и WebArena~\cite{zhou2023webarena}.

\paragraph{Риски самообучения.}
Самомодифицирующиеся агенты подвержены reward hacking~\cite{rewardhacking2025},
misevolution~\cite{misevolution2025} и иным hazards~\cite{selfimprovingrisks2025};
отсюда необходимость гейта верификации.

% =====================================================================
\section{Постановка задачи}
\label{sec:problem}

Рассматривается робот на арктической станции в частично известном
(или полностью неизвестном) мире.
Агент должен:
\begin{enumerate}[leftmargin=*]
  \item исследовать среду и строить пространственные убеждения;
  \item взаимодействовать с объектами (взять в «рюкзак», поставить,
        использовать~--- например, открыть дверь);
  \item решать целевые сценарии (сбор объектов, ревизия карты, реакция
        на нештатные ситуации);
  \item адаптироваться к динамике среды (снег, лёд, видимость, заряд).
\end{enumerate}

\textbf{Не} заявляется: физически точная термодинамика холода, деградация
сенсоров, sim-to-real, масштаб в тысячи сцен.

% =====================================================================
\section{Среда}
\label{sec:env}

\subsection{2D-симулятор и oracle-перцепция}

Основная экспериментальная площадка~--- 2D-симулятор.
Для каждого объекта хранится текстовое описание (внешний вид, надписи,
тип, разрешённые действия).
Текстовое описание сознательно заменяет модуль SLAM/восприятия и даёт
оракульное восприятие: агент получает безошибочную информацию об объектах
в радиусе видимости.

Наблюдения ограничены параметризуемым радиусом видимости (освещение~
--- модификатор).
Действия включают перемещение, поворот, взятие/установку предметов,
использование объектов и осмотр контейнеров
(контракт агент\(\leftrightarrow\)мир; детали API вне объёма статьи).

\subsection{Северная динамика}

\begin{itemize}[leftmargin=*]
  \item \textbf{Снег/заносы:} клетки становятся непроходимыми; объекты могут
        «теряться» в сугробе; между днями карта может оттаять.
  \item \textbf{Лёд:} параметрическая модель проскальзывания
        (вероятность и величина сдвига~--- функция скорости и типа поверхности),
        а не чистая «монетка».
  \item \textbf{Заряд:} аккумулятор разряжается; холодные клетки ускоряют расход;
        при низком заряде агент сам решает ехать заряжаться.
\end{itemize}

Робот узнаёт о свойствах клетки только проехав по ней; планировщик
перестраивает маршрут по мере пополнения памяти.

\subsection{Процедурная генерация карт}

Генератор даёт десятки вариантов карт по параметрам (плотность объектов,
доля льда/снега, размер, расположение станции/склада).
Фиксируются наборы train и \textbf{held-out} для честной оценки
метаагента~\cite{deitke2022procthor}.

\subsection{3D-визуализация}

Параллельно поддерживается 3D-симулятор на Unity с совпадающими картами~
--- для демонстрации («мультик») траекторий из лога 2D.
Связка не входит в основной научный результат.

% =====================================================================
\section{Метод}
\label{sec:method}

\subsection{LLM-агент-исследователь}

Агент работает в цикле ReAct~\cite{yao2022react}:
наблюдение \(\to\) рассуждение \(\to\) действие.
Ведётся текстовая память истории наблюдений; агент опирается на описания
объектов («колёса у гаража, лопата у склада»).
Рассматриваются две конфигурации моделей: облачные (через OpenRouter)
и локальные \(\le 30\mathrm{B}\).

\subsection{Абляция пространственной веры}

\begin{itemize}[leftmargin=*]
  \item \textbf{С координатами:} координаты доступны в памяти агента (упрощение).
  \item \textbf{Без координат:} объекты описываются относительно
        («дерево слева/справа»); агент строит и ревизует когнитивную карту.
        Опционально~--- belief-probing (выгрузка карты в JSON на шаге).
\end{itemize}
Сравнение режимов~--- прямая проверка Construct/Revise~\cite{theoryofspace2026}.

\subsection{Baseline'ы}

\begin{enumerate}[leftmargin=*]
  \item Детерминированная \textbf{спираль} (аналог «робот-пылесоса»): едет
        к однозначному ID, игнорирует описания.
  \item \textbf{Жадная/эвристическая} стратегия, читающая текстовые метки
        и приоритизирующая перспективные зоны.
\end{enumerate}

\subsection{Метаагент и гейт верификации}

Сильная модель анализирует логи роботов за день и предлагает улучшения:
правки промпта, новые скиллы, переписывание планировщика
(только класс планировщика с жёстким интерфейсом; физическое перемещение
недоступно для правок).

\textbf{Гейт верификации} (обязательный):
\begin{enumerate}[leftmargin=*]
  \item любая версия проходит регрессионные тесты на held-out картах;
  \item ухудшающие изменения автоматически откатываются;
  \item train и held-out жёстко разделены;
  \item мониторинг признаков reward hacking и неизменяемых safety-инвариантов.
\end{enumerate}

% =====================================================================
\section{Экспериментальная постановка}
\label{sec:setup}

\subsection{Сценарии}

\begin{itemize}[leftmargin=*]
  \item \textbf{«Собери все лопаты на склад»}~--- повтор по дням с вариацией
        расположения; автопроверка \(N\) из \(N\).
  \item \textbf{«Обновлённая карта»}~--- объект перемещён; агент отражает
        изменения (Belief Revision).
  \item \textbf{Реакция на новый объект} (например, «упавший человек»)~
        --- без захардкоженного поведения в промпте.
  \item \textbf{Взаимодействие с людьми}~--- «подслушивание» в радиусе
        слышимости и извлечение знания.
\end{itemize}

\subsection{Метрики}

SR (Success Rate), SPL~\cite{anderson2018spl}, покрытие карты / эффективность
исследования, число шагов/время, \(\mathrm{pass}^k\)~\cite{taubench2024}.
Автооценка~--- детерминированными правилами среды везде, где возможно;
ручная разметка~--- только для неоднозначных сценариев, с согласием аннотаторов.

\subsection{Сравнения и воспроизводимость}

\begin{itemize}[leftmargin=*]
  \item baseline'ы vs LLM; облачные vs локальные \(\le 30\mathrm{B}\);
  \item абляция «с/без координат»;
  \item эффект метаагента день-ко-дню на held-out;
  \item несколько сидов, доверительные интервалы;
  \item фиксация версий и дат API-моделей, сидов и гиперпараметров~\cite{pineau2018reproducibility}.
\end{itemize}

% =====================================================================
\section{Результаты}
\label{sec:results}

\todo{Заполнить после прогонов. Ниже~--- заготовки таблиц.}

\begin{table}[ht]
\centering
\caption{Сравнение агентов на held-out картах (заглушка)}
\label{tab:main}
\begin{tabular}{lcccc}
\toprule
Агент & SR \(\uparrow\) & SPL \(\uparrow\) & Покрытие \(\uparrow\) & \(\mathrm{pass}^k\) \(\uparrow\) \\
\midrule
Спираль & --- & --- & --- & --- \\
Эвристика & --- & --- & --- & --- \\
LLM (облако) & --- & --- & --- & --- \\
LLM (\(\le 30\mathrm{B}\)) & --- & --- & --- & --- \\
\bottomrule
\end{tabular}
\end{table}

\begin{table}[ht]
\centering
\caption{Абляция координат (заглушка)}
\label{tab:ablation}
\begin{tabular}{lccc}
\toprule
Режим & SR \(\uparrow\) & SPL \(\uparrow\) & Стабильность карты \\
\midrule
С координатами & --- & --- & --- \\
Без координат & --- & --- & --- \\
\bottomrule
\end{tabular}
\end{table}

\begin{table}[ht]
\centering
\caption{Эффект метаагента на held-out (заглушка)}
\label{tab:meta}
\begin{tabular}{lccc}
\toprule
День / итерация & SR & SPL & Принято гейтом \\
\midrule
0 (до самообучения) & --- & --- & --- \\
1 & --- & --- & --- \\
2 & --- & --- & --- \\
\bottomrule
\end{tabular}
\end{table}

% =====================================================================
\section{Обсуждение}
\label{sec:discussion}

Ожидаемые (проверяемые) гипотезы:
\begin{enumerate}[leftmargin=*]
  \item LLM-агент эффективнее или сопоставим с baseline'ами, но гибче
        на исследовании и сборе объектов.
  \item Северная динамика значимо влияет на поведение и выявляет слабые места
        (аналог Active-Passive Gap / Belief Inertia).
  \item Режим без координат показывает, умеет ли агент строить пространственную
        веру без готовой карты.
  \item Метаагент улучшает качество на held-out (а не только на train);
        порождённые скиллы ускоряют похожие новые задачи.
  \item Локальные \(\le 30\mathrm{B}\) уступают frontier-моделям в многошаговом
        планировании~\cite{liu2023agentbench}; это результат, а не неудача;
        метаагент может частично «подтягивать» слабых исполнителей.
\end{enumerate}

\todo{После результатов~--- обсудить расхождения с гипотезами и случаи
reward hacking / откатов гейта.}

% =====================================================================
\section{Ограничения}
\label{sec:limitations}

\begin{enumerate}[leftmargin=*]
  \item \textbf{Oracle-перцепция.} Текстовые описания заменяют реальное
        восприятие/SLAM. Результаты~--- верхняя граница планирования при
        идеальном восприятии.
  \item \textbf{Отсутствие физики холода.} Энергетика, термодинамика,
        деградация сенсоров не моделируются; «северность»~--- семантика
        и динамика задач.
  \item \textbf{Только симуляция.} Sim-to-real не заявляется и не проверяется.
  \item \textbf{Ограниченный масштаб.} Десятки процедурно сгенерированных карт~
        --- достаточно для статистики, но меньше индустриальных бенчмарков.
  \item \textbf{Упрощённая стохастика скольжения}~--- не физика контакта.
  \item \textbf{Самообучение под контролем гейта} снижает, но не устраняет
        риск reward hacking; наблюдения ограничены длиной эксперимента.
\end{enumerate}

% =====================================================================
\section{Заключение}
\label{sec:conclusion}

Мы описываем лёгкий прототип для LLM-агентов активного исследования
частично известного северного мира с oracle-перцепцией, абляцией
пространственной веры и метаагентом под гейтом верификации.
Цель~--- честно измерить верхнюю границу планирования и надёжность
самоулучшения в симуляции с явной северной семантикой задач.

\todo{Кратко суммировать итоговые числа и открытые вопросы после экспериментов.}

\printbibliography

\end{document}
